\scanpage{136}
\mspage{V.8}
% archive sheet number: 135
% manuscript lines 1-4
extension \inserted{algébrique} séparable de $K(T_1,\ldots,T_r)$, donc une ext. séparable de $K=k(\unsure{\eta|y})$, (et $\lambda=\operatorname{Spec}L\simeq(S_1)_\eta$ \deleted{\unsure{est un}}). D'après 2.1.1. appliqué à la situation $X_1\to S_1$, et $U_1\subset X_1$, $X'\mathbin{\unsure{-|\to}}X_1$, on voit que
\[
H^i(U'\times_{S_1}\bar\lambda,\Lambda)=0\quad\text{si }i>n-r,\qquad
\text{i.e. }R^i\varphi_*(\Lambda_{U'_{\bar\eta}})=0\quad\text{si }i>n-r
\tag*{(2.\unsure{2|1}.1)}
\]
% manuscript lines 5-8
D'autre part, \deleted{$\bar\lambda$ \unsure{est un}} on sait que \deleted{\unsure{$S_1$}} ($S_1\mathbin{\unsure{-|\to}}S$ étant déduit par localisation stricte du morphisme lisse $Y\to S$) que $S_{1\eta}=\lambda$ est acyclique pour les nos premiers $\ne p$ (SGA 4 XV). Donc la suite spectrale
% slanted
\margin{réf. inutile, on voit directement que $\operatorname{Gal}(\unsure{\bar\lambda|\bar\eta}/S_{1\eta})$ est un pro-$p$-groupe\unsure{---}.}
\[
H^i(U'_{\bar\eta},\Lambda)\Longleftarrow\mathbb E_{\unsure{2|,}}^{p,q}=H^p(\lambda,R^q\varphi_*(\Lambda_{U'_{\bar\eta}}))
\]
% manuscript lines 9-13
dégénère et donne
\[
H^i(U'_{\bar\eta},\Lambda)\simeq H^0(\lambda,R^i\varphi_*(\Lambda_{U'_{\bar\eta}}))\hookrightarrow H^i(U'\times_{S_1}\bar\lambda,\Lambda)
\tag*{(2.2.2)}
\]
ce qui implique par (2.2.1)
\[
H^i(U'_{\bar\eta},\Lambda)=0\qquad\text{pour }i>n-r,
\]
et achève la démonstration. On a prouvé :

% manuscript lines 14-17
\uline{Théorème} 2.3. Avec les notations de 2.1, on a \inserted{pour $x\in X_0$,}
\[
(\Psi^q_{\unsure{\nu|v}})_{\bar x}=0\qquad\text{pour }q>n_x-\dim\overline{\{x\}}
\]
(et à fortiori pour $i>\dim\underline{\unsure{\mathcal O|O}}_{X_0,x}=\dim_x X_{\unsure{x|0}}-\dim\overline{\{x\}}$).
De façon équivalente on a
\[
(\underline\Phi^q_\nu)_{\bar x}=0\quad\text{si }q>n_x+1-\dim\overline{\{x\}},\quad q\geq1.
\]
% manuscript lines 18-23: cancelled block
\deleted{Appliquant le théorème de dim. cohomologique SGA 4 \uline{IX}, on trouve}
\deleted{\inserted{si $n=\dim X_\eta$, $q\geq1$, $X=\overline{X_\eta}$ donc $\underline\Phi^0_\nu=0$ :}}
\[
\deleted{H^p(X_0,(\underline\Phi^q_\nu)|X_0)=0\quad\text{si }p\geq2(n+1-q),}
\]
\deleted{i.e. $p+q\geq2n+2-q$}

\deleted{\uline{Corollaire} 2.\unsure{3}. On a, si $X$ propre sur $S$}

% manuscript lines 24-26
\uline{Corollaire} 2.\unsure{4|3}. On a : $\underline\Psi^q_\nu=0$ si $q>\dim X_\eta$, $\underline\Phi^q_\nu=0$ si $q>\dim X_\eta+1$.
